\newpage
\section*{Solucionario completo}

\noindent Cada respuesta incluye la clave y una justificación breve. En los ejercicios numéricos se indica la ecuación o el cálculo principal.

\subsection*{Bloque A. Fundamentos teóricos y conceptuales}

\sol{1}{C}{La competencia perfecta supone muchos compradores y vendedores, producto homogéneo, información perfecta, libre entrada y salida, y empresas precio-aceptantes.}

\sol{2}{A}{La empresa individual enfrenta una demanda horizontal: puede vender cualquier cantidad al precio de mercado, pero no puede fijar un precio superior.}

\sol{3}{C}{La demanda del mercado es descendente: a mayor precio, menor cantidad demandada por el conjunto de consumidores.}

\sol{4}{D}{La regla de maximización es $IMg=CMg$. En competencia perfecta, $IMg=P$, por lo que también $P=CMg$.}

\sol{5}{C}{Como la empresa vende cada unidad adicional al mismo precio, $IMg=P$.}

\sol{6}{D}{El beneficio económico es $IT-CT$ económico, donde $CT$ incluye costos explícitos e implícitos o de oportunidad.}

\sol{7}{B}{El beneficio normal remunera el costo de oportunidad del capital propio. Cuando el beneficio económico es cero, la empresa obtiene exactamente su beneficio normal.}

\sol{8}{D}{En el corto plazo, el beneficio puede ser positivo si $P>CMe$, cero si $P=CMe$, o negativo si $P<CMe$, siempre que decida producir.}

\sol{9}{B}{Si $P<CVMe$, la empresa no cubre ni los costos variables; le conviene suspender. Si $CVMe<P<CMe$, cubre variables y parte de los fijos.}

\sol{10}{D}{La oferta de corto plazo es el tramo de $CMg$ situado por encima del mínimo $CVMe$; por debajo de ese precio, la empresa no produce.}

\sol{11}{B}{La oferta de la industria es la suma horizontal de las cantidades ofrecidas por cada empresa a cada precio.}

\sol{12}{C}{Con libre entrada y salida, el equilibrio de largo plazo implica $P=CMg=CMe$ y beneficio económico nulo; las empresas solo obtienen beneficio normal.}

\sol{13}{D}{Si $P>\min CMe$, hay beneficios extraordinarios, entran empresas, la oferta aumenta y el precio cae hasta el mínimo del costo medio.}

\sol{14}{A}{Si $P<\min CMe$, las empresas incurren en pérdidas, salen del mercado, la oferta disminuye y el precio sube.}

\sol{15}{C}{En una industria de costo constante, los precios de los insumos no cambian al expandirse la industria; la oferta de largo plazo es horizontal.}

\sol{16}{A}{Si al expandirse la industria los costos de los insumos aumentan, la oferta de largo plazo tiene pendiente positiva.}

\sol{17}{D}{Si la expansión reduce los costos de los insumos, la oferta de largo plazo tiene pendiente negativa.}

\sol{18}{C}{El excedente del productor es $IT-CV$. También puede medirse como el área situada entre el precio y la curva de costo marginal.}

\sol{19}{A}{El excedente del consumidor es la diferencia entre la disposición a pagar y el precio efectivamente pagado.}

\sol{20}{A}{El excedente total es $EC+EP$ y se maximiza en el equilibrio competitivo, donde $P=CMg$.}

\sol{21}{C}{Un precio máximo por debajo del equilibrio reduce la cantidad ofrecida y aumenta la demandada: genera escasez y pérdida de bienestar.}

\sol{22}{A}{Un precio mínimo por encima del equilibrio genera excedente de oferta y pérdida de bienestar.}

\sol{23}{A}{La carga económica no depende solo de quién paga legalmente el impuesto, sino de las elasticidades: el lado menos elástico soporta una mayor proporción.}

\sol{24}{B}{En un subsidio, el beneficio se reparte según las elasticidades; el lado menos elástico captura una mayor parte.}

\sol{25}{C}{La eficiencia asignativa requiere que el precio refleje el costo marginal: $P=CMg$.}

\sol{26}{C}{En el largo plazo competitivo, la libre entrada lleva a producir en el mínimo del costo medio: $P=\min CMe$.}

\sol{27}{D}{La libre entrada elimina los beneficios extraordinarios: si hay beneficios, entran empresas hasta que $P=\min CMe$ y el beneficio económico se anula.}

\sol{28}{D}{Una empresa precio-aceptante no puede influir en el precio de mercado; lo toma como dato y decide solo su cantidad.}

\sol{29}{A}{El equilibrio de mercado se alcanza cuando la cantidad demandada iguala la cantidad ofrecida: $Q_d=Q_s$.}

\sol{30}{A}{En el corto plazo, la empresa produce donde $P=CMg$, siempre que el precio cubra al menos el costo variable medio.}

\subsection*{Bloque B. Aplicaciones cuantitativas y optimización}

\sol{31}{B}{$CMg=10+2q$. Igualando $P=30$: $10+2q=30$, de donde $q=10$.}

\sol{32}{A}{$IT=30(10)=300$. $CT=100+10(10)+10^2=300$. Beneficio $=300-300=0$.}

\sol{33}{C}{$CMg=5+q$. Igualando $P=25$: $5+q=25$, de donde $q=20$.}

\sol{34}{B}{$IT=25(20)=500$. $CT=50+5(20)+0.5(20)^2=50+100+200=350$. Beneficio $=150$.}

\sol{35}{A}{$CMg=4+2q=20$, entonces $q=8$. $IT=160$; $CT=100+4(8)+8^2=196$; beneficio $=-36$.}

\sol{36}{D}{$CVMe=4+q$, que se minimiza en $q=0$ con valor 4. Si $P<4$, la empresa no cubre los costos variables y suspende.}

\sol{37}{B}{$CMg=2q+10$. Igualando $P=30$: $2q+10=30$, de donde $q=10$.}

\sol{38}{A}{La oferta individual es $q=(P-10)/2$. Con 100 empresas: $Q=100(P-10)/2=50P-500$.}

\sol{39}{A}{$1000-50P=50P-500$, entonces $1500=100P$ y $P=15$.}

\sol{40}{D}{Sustituyendo $P=15$: $Q=1000-50(15)=250$.}

\sol{41}{B}{$q=(P-10)/2=(15-10)/2=2.5$. También $250/100=2.5$.}

\sol{42}{C}{$CMe=q+4+100/q$. Derivando: $1-100/q^2=0$, de donde $q=10$ y $CMe=10+4+10=24$. En largo plazo, $P=24$.}

\sol{43}{A}{El mínimo del costo medio se alcanza en $q=10$.}

\sol{44}{C}{Con $P=24$: $Q=500-10(24)=260$. Si cada empresa produce 10, el número de empresas es $260/10=26$.}

\sol{45}{B}{La demanda inversa es $P=50-0.1Q$. $EC=0.5(50-24)(260)=3380$.}

\sol{46}{C}{En el equilibrio de largo plazo con libre entrada, el beneficio económico de cada empresa es cero.}

\sol{47}{A}{En una industria de costo constante, la oferta de largo plazo es horizontal al mínimo del costo medio. Un aumento de demanda eleva la cantidad, pero el precio vuelve a 24.}

\sol{48}{B}{$120-2P=3P-30$, entonces $150=5P$ y $P=30$.}

\sol{49}{C}{Sustituyendo $P=30$: $Q=120-2(30)=60$.}

\sol{50}{D}{La demanda inversa es $P=60-0.5Q$. $EC=0.5(60-30)(60)=900$.}

\sol{51}{C}{La oferta inversa es $P=10+Q/3$. $EP=0.5(30-10)(60)=600$.}

\sol{52}{D}{$ET=EC+EP=900+600=1500$.}

\sol{53}{D}{Con $P=20$: $Q_d=120-40=80$ y $Q_s=60-30=30$. La escasez es $80-30=50$.}

\sol{54}{A}{En $Q=30$, el precio de demanda es $(120-30)/2=45$. La pérdida es $0.5(80-30)(45-20)=625$.}

\sol{55}{D}{Con impuesto de 5, la oferta es $Q_s=3(P-5)-30=3P-45$. Igualando: $120-2P=3P-45$, de donde $P=33$.}

\sol{56}{B}{Con $P=33$: $Q=120-2(33)=54$.}

\sol{57}{A}{Recaudación $=5(54)=270$.}

\sol{58}{D}{$PCP=0.5(5)(60-54)=15$.}

\sol{59}{C}{El precio sube de 30 a 33: los consumidores soportan 3 por unidad.}

\sol{60}{B}{El productor recibe $33-5=28$, frente a 30 inicial: soporta 2 por unidad.}

\sol{61}{A}{Con subsidio de 5, la oferta es $Q_s=3(P+5)-30=3P-15$. Igualando: $120-2P=3P-15$, de donde $P=27$.}

\sol{62}{A}{Con $P=27$: $Q=120-2(27)=66$. Costo del subsidio $=5(66)=330$.}

\sol{63}{C}{$PCP=0.5(5)(66-60)=15$.}

\sol{64}{B}{En una industria de costo constante, la oferta de largo plazo es horizontal.}

\sol{65}{C}{En una industria de costo creciente, la oferta de largo plazo es ascendente.}

\subsection*{Bloque C. Casos con implicancias tarifarias y regulatorias}

\sol{66}{B}{$80-2P=2P-20$, entonces $100=4P$ y $P=25$.}

\sol{67}{B}{La demanda inversa es $P=40-0.5Q$. $EC=0.5(40-25)(30)=225$.}

\sol{68}{D}{La oferta inversa es $P=10+0.5Q$. $EP=0.5(25-10)(30)=225$.}

\sol{69}{C}{$ET=225+225=450$.}

\sol{70}{B}{Con $P=20$: $Q_d=80-40=40$ y $Q_s=40-20=20$. La escasez es 20.}

\sol{71}{D}{En $Q=20$, el precio de demanda es $(80-20)/2=30$. $PCP=0.5(40-20)(30-20)=100$.}

\sol{72}{B}{Con impuesto de 4, la oferta es $Q_s=2(P-4)-20=2P-28$. Igualando: $80-2P=2P-28$, de donde $P=27$.}

\sol{73}{D}{Con $P=27$: $Q=80-2(27)=26$. Recaudación $=4(26)=104$.}

\sol{74}{A}{$PCP=0.5(4)(30-26)=8$.}

\sol{75}{D}{El precio sube de 25 a 27: los consumidores soportan 2 por unidad.}

\sol{76}{B}{El productor recibe $27-4=23$, frente a 25 inicial: soporta 2 por unidad.}

\sol{77}{C}{Con subsidio de 4, la oferta es $Q_s=2(P+4)-20=2P-12$. Igualando: $80-2P=2P-12$, de donde $P=23$.}

\sol{78}{C}{Con $P=23$: $Q=80-2(23)=34$. Costo del subsidio $=4(34)=136$.}

\sol{79}{A}{$PCP=0.5(4)(34-30)=8$.}

\sol{80}{B}{Si $P>\min CMe$, hay beneficios extraordinarios; entran empresas, la oferta aumenta y el precio tiende a bajar hasta el mínimo del costo medio.}

\sol{81}{A}{En una industria de costo constante, la oferta de largo plazo es horizontal.}

\sol{82}{D}{En un monopolio natural, $P=CMg<CMe$ genera pérdidas; se requiere un subsidio o un cargo fijo para cubrirlas.}

\sol{83}{B}{Un mercado contestable puede aproximarse al resultado competitivo si la entrada potencial es creíble, aunque haya pocas empresas.}

\sol{84}{B}{Una reducción de la tarifa aumenta el excedente del consumidor, pues paga menos por la misma cantidad o consume más.}

\sol{85}{C}{Una reducción de la tarifa disminuye el excedente del productor y puede comprometer la recuperación de costos y la sostenibilidad financiera.}

\sol{86}{B}{Un mercado de derechos de agua permite que el recurso se reasigne hacia los usos de mayor valor económico, sujeto a los límites ambientales.}

\sol{87}{B}{El impuesto pigouviano óptimo iguala el costo externo marginal, internalizando la externalidad negativa.}

\sol{88}{C}{Un precio máximo por debajo del equilibrio genera escasez; la asignación se realiza por mecanismos no basados en el precio, como racionamiento o colas.}

\sol{89}{D}{Un subsidio directo focalizado protege el ingreso de los hogares vulnerables sin distorsionar el precio de mercado ni generar escasez.}

\sol{90}{D}{El equilibrio competitivo es eficiente en sentido asignativo, pero no garantiza equidad; pueden requerirse transferencias o subsidios focalizados.}
